\section*{Capítulo \ref{ch:g4:numbers} --- Números grandes}

\begin{solution}{exo:g4:numbers:1}
$3\,520\,000$; \quad $7\,000\,008$; \quad $240\,060$.
\end{solution}

\begin{solution}{exo:g4:numbers:2}
$5\,204\,300$: cinco millones doscientos cuatro mil trescientos.

$10\,050\,000$: diez millones cincuenta mil.

$999\,999$: novecientos noventa y nueve mil novecientos
noventa y nueve.
\end{solution}

\begin{solution}{exo:g4:numbers:3}
Decenas de miles: $7$. Clase de cientos de millones: no hay
dígitos de cientos de millones en un número de siete dígitos --- los millones
clase de $8\,472\,156$ contiene solo el dígito $8$ (unidades de
millones). Miles contenidos: $8\,472$ (desde
$8\,472\,156 = 8\,472 \times 1\,000 + 156$).
\end{solution}

\begin{solution}{exo:g4:numbers:4}
$980\,000 < 1\,020\,000$; \quad
$3\,456\,000 < 3\,465\,000$ (clase de miles $456 < 465$); \quad
$7\,090\,050 < 7\,090\,500$.
\end{solution}

\begin{solution}{exo:g4:numbers:5}
$980\,500 < 2\,030\,000 < 2\,299\,999 < 2\,300\,000$.
\end{solution}

\begin{solution}{exo:g4:numbers:6}
A los diez: $63\,780$. A los cien: $63\,800$. A los mil:
$64\,000$. A los diez mil: $60\,000$.
\end{solution}

\begin{solution}{exo:g4:numbers:7}
$5\,000\,000$; \quad $6\,000\,000$ (el $499\,999$ no alcanza
\omterm{def:g3:division:half}{medio} un millón); \quad $10\,000\,000$ (exactamente a mitad de camino rondas arriba).
\end{solution}

\begin{solution}{exo:g4:numbers:8}
$998\,500 \to 999\,500 \to 1\,000\,500 \to 1\,001\,500$: el segundo
El paso cruza el millón y se ilumina una nueva clase a la izquierda.

$950\,000 \to 1\,050\,000 \to 1\,150\,000$: mismo cruce en el
primer paso.
\end{solution}

\begin{solution}{exo:g4:numbers:9}
Ambos tienen razón. ``Más de $80\,000$'' rondas hasta las diez
mil ($81\,338 \to 80\,000$); ``sobre $81\,000$'' rondas al
mil más cercano ($81\,338 \to 81\,000$).
\end{solution}

\begin{solution}{exo:g4:numbers:10}
Más grande: $7\,654\,321$. Pequeñísimo: $1\,234\,567$. Más cercano a
$4\,000\,000$: justo debajo funciona mejor --- $3\,765\,421$ es
$234\,579$ lejos, mientras que justo arriba, $4\,123\,567$, es solo
$123\,567$ de distancia: lo más cercano es $4\,123\,567$.
\end{solution}

\begin{solution}{exo:g4:numbers:11}
redondeo al millar más cercano da $50\,000$ para todos los números
de $49\,500$ (a \omterm{def:g3:division:half}{medio} camino rondas arriba a $50\,499$. Más grande:
$50\,499$; pequeñísimo: $49\,500$.
\end{solution}
